\documentclass[11pt]{article}
\usepackage[T1]{fontenc}
\usepackage{amsmath,amssymb,amsthm}
\usepackage[margin=1in]{geometry}
\usepackage[hidelinks]{hyperref}

\newtheorem{theorem}{Theorem}
\newtheorem{lemma}[theorem]{Lemma}

\newcommand{\PsiHG}{\Psi_{H,g}}
\newcommand{\coeff}[2]{[#1]#2}

\title{Loopy-Polynomial Injectivity for Rooted Trees\\of Height at Most Two}
\author{Xiaojian Zeng\\
PBC School of Finance, Tsinghua University, 100190\\
\texttt{xiaojianzeng@foxmail.com}}
\date{}

\begin{document}
\maketitle

\begin{abstract}
For an integer $g\geq 3$, a specialization of the rooted
$U$-polynomial records the order of the root component by replacing its
marker with the ordinary variable indexed by $g-1$ plus that order.  We
prove that this specialization distinguishes all rooted trees of root
height at most two, for every order and without a restriction on branch
sizes.  The proof extracts the elementary symmetric functions of the
numbers of depth-two leaves in the root branches, together with the root
degree, from two explicit families of coefficients.  This resolves only
the height-at-most-two specialization of the root-injectivity conjecture;
the conjecture for unrestricted rooted trees is not addressed.
\end{abstract}

\section{Introduction}

Let $(H,\rho)$ be a finite rooted tree and let $g\geq 3$.  For
$J\subseteq E(H)$, write $c(J)$ for the components of the spanning forest
$(V(H),J)$ and $R_J$ for the component containing $\rho$.  Define
\begin{equation}\label{eq:def}
 \PsiHG(\mathbf z)=
 \sum_{J\subseteq E(H)}
 z_{g-1+|V(R_J)|}
 \prod_{\substack{D\in c(J)\\D\neq R_J}}z_{|V(D)|},
\end{equation}
where $z_1,z_2,\ldots$ are commuting indeterminates.

Kirillov, Nenashev, Shapiro, and Vaintrob introduced this polynomial in a
reduction for unicyclic graphs and conjectured that, for every $g\geq3$,
it distinguishes all rooted trees of each fixed order
\cite[Conjecture~9.10]{KirillovEtAl2026}.  They verified the conjecture
through order $13$, observed that the range $g\geq |V(H)|$ is automatic,
and obtained an unconditional result when the relevant branches are
smaller than $g$ \cite[Remark~9.12 and Theorem~9.14]{KirillovEtAl2026}.
The fully pointed precursor records the root-component order in a separate
variable.  It is the rooted $U$-polynomial of Aliste-Prieto, de Mier, and
Zamora, which reconstructs every rooted tree
\cite[Theorem~9]{AlistePrietoDeMierZamora2021}.

Component-order data also appear in reconstruction from other graph
polynomials.  For example, Zeng uses leading star-basis coefficient layers
to recover weighted cores under a multiplicity-isolation condition and
gives an integer-partition encoding of height-two rooted branches in a
diameter-six setting \cite[Sections~2--4]{Zeng2026Multiplicity}.  The
argument here is independent of those results: the height-two normal form
allows the coefficients of \eqref{eq:def} itself to recover the complete
branch multiset, with repetitions allowed.

Our contribution is the following exact restricted case of the cited
root-injectivity conjecture.  It removes every branch-size restriction
within the class of root height at most two.  It makes no assertion for
larger root height and no claim of priority beyond the bounded comparison
above.

\begin{theorem}\label{thm:main}
Let $g\geq3$ and $N\geq1$.  If $(H,\rho)$ and $(H',\rho')$ are rooted
trees on $N$ vertices, both of root height at most two, then
\[
 \Psi_{H,g}=\Psi_{H',g}
 \quad\Longrightarrow\quad
 (H,\rho)\cong(H',\rho').
\]
\end{theorem}

\section{Coefficient extraction and reconstruction}

For $N\geq2$, deleting $\rho$ leaves $r=\deg_H(\rho)$ components.
The height assumption and the fact that $H$ is a tree imply that each is
a star centered at its unique neighbor of $\rho$.  Write their orders as
$b_1,\ldots,b_r$, and put
\begin{equation}\label{eq:normal}
 d_i=b_i-1\geq0,
 \qquad \sum_{i=1}^r b_i=N-1.
\end{equation}
Thus $d_i$ is the number of depth-two leaves in branch $i$.  Conversely,
the multiset $\{\!\{b_1,\ldots,b_r\}\!\}$ determines the rooted
isomorphism type.

\begin{lemma}\label{lem:coefficients}
For $0\leq k\leq\lfloor(N-1)/2\rfloor$, let
\[
 A_k=\coeff{z_gz_2^kz_1^{N-1-2k}}{\PsiHG}.
\]
Then
\begin{equation}\label{eq:Ak}
 A_k=e_k(d_1,\ldots,d_r).
\end{equation}
Moreover,
\begin{equation}\label{eq:B}
 \coeff{z_{g+1}z_1^{N-2}}{\PsiHG}=r.
\end{equation}
\end{lemma}

\begin{proof}
In every summand of \eqref{eq:def}, the root-component factor has index
at least $g$.  The target monomial defining $A_k$ has exactly one factor
whose index is at least $g$, namely $z_g$; here the hypothesis $g\geq3$
also covers the endpoint $g=3$.  Thus the root component is a singleton,
so every root edge is absent, and the factor $z_g$ has already been
accounted for.

Consider the detached star in branch $i$.  If $q$ of its $d_i$ leaf
edges are selected, its center lies in a component of order $q+1$ and
each unselected leaf is isolated.  The remaining target factors have
indices only $1$ and $2$, so $q$ is either zero or one.  Selecting one
leaf edge in branch $i$ can be done in $d_i$ ways.  A contributing forest
is therefore specified uniquely by choosing $k$ distinct branches and
one leaf edge in each, giving
\[
 A_k=\sum_{\substack{S\subseteq\{1,\ldots,r\}\\|S|=k}}
      \prod_{i\in S}d_i=e_k(d_1,\ldots,d_r).
\]
This includes $k=0$, for which the empty forest is the unique
contributor.  It also excludes a hidden commutative-monomial collision:
an additional component of order $g$ would contribute a second $z_g$,
and any order outside $\{1,2\}$ would contribute a factor absent from the
target.

For \eqref{eq:B}, the target's only factor with index at least $g$ is
$z_{g+1}$.  Hence the root component has order two.  Every other vertex
is isolated, so the selected edge set consists of exactly one of the $r$
edges incident with the root.  Conversely, each such one-edge set gives
the target monomial.  Its coefficient is therefore $r$.
\end{proof}

\begin{proof}[Proof of Theorem~\ref{thm:main}]
The case $N=1$ is immediate: there is one rooted tree and its polynomial
is $z_g$.  Suppose $N\geq2$ and set
$L=\lfloor(N-1)/2\rfloor$.  If
\[
 m=|\{i:d_i>0\}|,
\]
then every branch counted by $m$ has at least two vertices.  Hence
$2m\leq N-1$, so $m\leq L$.  Lemma~\ref{lem:coefficients} consequently
recovers the entire polynomial
\begin{equation}\label{eq:Q}
 Q_H(t)=\sum_{k=0}^{L}A_kt^k
       =\prod_{i=1}^{r}(1+d_it).
\end{equation}
The multiset of positive $d_i$ is determined, with multiplicity, by the
roots $-1/d_i$ of $Q_H$ (equivalently, by unique factorization of its
normalized nonconstant linear factors).  Equation \eqref{eq:B} determines
$r$, so appending exactly $r-m$ zeros recovers the full multiset of the
$d_i$.  Adding one to every entry recovers the branch-order multiset in
\eqref{eq:normal}, and hence the rooted tree.

If $\Psi_{H,g}=\Psi_{H',g}$ for the same $g$ and $N$, all coefficients
used above agree.  Thus the two branch-order multisets agree and the
rooted trees are isomorphic.
\end{proof}

\section{Exact finite audit}

An ancillary Python program independently checks the coefficient formulas
and injectivity over a finite range.  For every $1\leq N\leq14$, it
enumerates every integer partition of $N-1$, constructs the corresponding
height-two rooted tree, enumerates all $2^{N-1}$ edge subsets, and computes
\eqref{eq:def} from the resulting forest components using exact integer
arithmetic.  For every $3\leq g\leq14$, all tested instances satisfy
\eqref{eq:Ak}, \eqref{eq:B}, and full-polynomial injectivity.  From the
project root the committed output is reproduced by
\begin{verbatim}
python3 evidence/audit_height_two.py --max-n 14 --max-g 14
\end{verbatim}
This audit is finite corroboration only.  No computational output is used
to infer any untested case; the quantified theorem rests on the proof
above.

\section*{Computational and AI assistance disclosure}

OpenAI Codex language models assisted with mathematical exploration,
proof drafting, exact-program development, and checks of the argument and
manuscript.  A fresh model context reconstructed the proof from the
definition and reran the committed finite audit.  The author reviewed the
result and takes responsibility for the claims.  Neither model output nor
finite computation is a premise of Theorem~\ref{thm:main}.

\section*{Data and code availability}

The exact Python verifier and its committed JSON output accompany this
article in the evidence directory, under the filenames
\texttt{audit\_height\_two.py} and \texttt{audit\_height\_two.json}.

\bibliographystyle{plain}
\bibliography{references}

\end{document}
